\documentclass[11pt,a4paper]{article}
\usepackage[T1]{fontenc}
\usepackage[utf8]{inputenc}
\usepackage{amsmath,amssymb,amsthm}
\usepackage[margin=1in]{geometry}
\usepackage[colorlinks=true,linkcolor=blue,urlcolor=blue]{hyperref}
\theoremstyle{plain}
\newtheorem{theorem}{Theorem}
\newtheorem{lemma}[theorem]{Lemma}
\newtheorem{proposition}[theorem]{Proposition}
\newtheorem{corollary}[theorem]{Corollary}
\theoremstyle{definition}
\newtheorem{remark}[theorem]{Remark}

\title{The empirical recurrence for OEIS A183716, proved by transfer matrix}
\author{Adrian Perez Fontelles\\ \small Independent researcher}
\date{3 September 2026}
\begin{document}
\maketitle

\begin{abstract}
OEIS A183716 counts arrays of width $7$ over $\{0,1,2,3,4\}$ subject to a condition read round the
perimeter of every $2\times2$ subblock: every subblock turns strictly one way with a single decrease. The entry carries an empirical recurrence of
order $35$ contributed by R.~H.~Hardin. It is true, and it is decidable rather than
empirical: the condition is settled inside a bounded window of consecutive rows, so the
count is a walk count on a finite digraph with $S=78126$ vertices, and whether a linear
recurrence holds for such a count is a finite computation.
\end{abstract}

\noindent\small 2020 Mathematics Subject Classification. 05A15, 05A16, 15A18.\normalsize

\section{The conjecture}

OEIS A183716 is ``1/20 of the number of (n+1) X 7 0..4 arrays with every 2 X 2 subblock strictly increasing clockwise or counterclockwise with one decrease.''. It has offset $1$ and begins
\[
169, 1797, 16964, 169560, 1683197, 16823812, 168070828, 1680902120,\ \dots
\]
The entry states, as a conjecture contributed by its author and never marked settled:
\begin{quote}
Empirical: a(n)=12*a(n-1)+41*a(n-2)-731*a(n-3)-236*a(n-4)+17377*a(n-5)-11599*a(n-6)-210068*a(n-7)+239143*a(n-8)+1435333*a(n-9)-2045028*a(n-10)-5829087*a(n-11)+9632462*a(n-12)+14361840*a(n-13)-27309714*a(n-14)-21401634*a(n-15)+48701640*a(n-16)+18436650*a(n-17)-56085498*a(n-18)-7403128*a(n-19)+42308694*a(n-20)-1027934*a(n-21)-20976128*a(n-22)+2642614*a(n-23)+6772235*a(n-24)-1305140*a(n-25)-1387863*a(n-26)+326941*a(n-27)+171884*a(n-28)-44771*a(n-29)-11807*a(n-30)+3188*a(n-31)+387*a(n-32)-103*a(n-33)-4*a(n-34)+a(n-35).
\end{quote}
As of the ``Last modified'' line on the live entry (Oct 06 2015 21:59:20, revision 6) this is still
recorded as empirical, and nothing on the entry records it as proved.

\section{What is being counted}

The arrays have $7$ columns and a growing number of rows, with entries in $\{0,1,2,3,4\}$. A
$2\times2$ subblock is a set of four cells
\[
B=\begin{pmatrix} p & q \\ r & s \end{pmatrix},
\qquad p=x(i,j),\; q=x(i,j+1),\; r=x(i+1,j),\; s=x(i+1,j+1).
\]
Its perimeter is a cycle of four cells, walked clockwise as $p\to q\to s\to r\to p$ and
counterclockwise as $p\to r\to s\to q\to p$. Write
\[
\mathrm{cw}(B)=[p<q]+[q<s]+[s<r]+[r<p],\qquad
\mathrm{ccw}(B)=[p<r]+[r<s]+[s<q]+[q<p]
\]
for the number of steps at which the value rises, going each way round.

\begin{lemma}\label{lem:sum}
$\mathrm{cw}(B)+\mathrm{ccw}(B)=4-\#\{\text{edges of the cycle with equal ends}\}$.
\end{lemma}

\begin{proof}
Each of the four edges of the square is walked once by each cycle, and in opposite directions.
An edge with distinct ends therefore contributes exactly one rise, to one of the two counts;
an edge with equal ends contributes to neither.
\end{proof}

The condition the entry imposes is this: every subblock $B$ has its four corners pairwise unequal along the cycle ($p\ne q$, $q\ne s$, $s\ne r$, $r\ne p$) and
\[
\mathrm{cw}(B)\in\{1,3\}.
\]
Going round clockwise there are then four strict steps, of which either three rise and one falls ($\mathrm{cw}(B)=3$: strictly increasing clockwise with one decrease) or one rises and three fall ($\mathrm{cw}(B)=1$, which is the same statement read counterclockwise). Call the subblock \emph{clockwise} in the first case and \emph{counterclockwise} in the second; this is its increasing direction. The entry publishes $1/20$ of that count; the division is by the constant 20 and changes nothing about the recurrence.

\section{A bounded window}

Two consecutive rows decide every subblock lying between them, and the condition of Section 2 is a condition on single subblocks only. So the state is the row just written, and nothing else.

One step of the walk writes a whole row. The row is filled left to right, and the subblock
closed by positions $j-1,j$ is tested the moment position $j$ is written, so a partial row that
cannot be completed is abandoned rather than enumerated. The walk starts at the state standing
for the empty array, so that a walk of $n$ steps is an array of $n$ rows.

The end vector is all ones: an array may stop after any row. Thus
\[
a(n)\;=\;\iota^{\!\top}M^{\,j}\tau ,\qquad S=78126,
\]
for the transfer matrix $M$ on those $S$ states. The walk index $j$ and the entry's own $n$
differ by a constant, read off by matching the published terms rather than assumed. In
particular $a$ is $C$-finite, and every question about linear recurrences it satisfies is
decidable.

\section{The proof}

\begin{lemma}\label{lem:crit}
Let $q(t)=t^{35}-\sum_i c_it^{35-i}$ be the characteristic polynomial of the
conjectured recurrence. The residual $a(n)-\sum_ic_ia(n-i)$ equals
$\iota^{\!\top}M^{\,j}q(M)\tau$ for the corresponding walk index $j$, so the recurrence
holds from some point on if and only if $\iota^{\!\top}M^{j}q(M)\tau=0$ for all large $j$,
and it is enough to examine $j<S$.
\end{lemma}

\begin{proof}
Factor $M^{j}$ out of $M^{j+35}-\sum_ic_iM^{j+35-i}$. For the bound, $M^{S}$
is an integer combination of $I,M,\dots,M^{S-1}$ by Cayley--Hamilton, so the numbers
$u_j=\iota^{\!\top}M^{j}q(M)\tau$ obey the monic recurrence given by the characteristic
polynomial of $M$; $S$ consecutive zeros force every later one, and the last nonzero $u_j$
pins the threshold exactly.
\end{proof}

\begin{theorem}
\[
a(n)\;=\;12\,a(n-1) + 41\,a(n-2) - 731\,a(n-3) - 236\,a(n-4) + 17377\,a(n-5) - 11599\,a(n-6) - 210068\,a(n-7) + 239143\,a(n-8) + 1435333\,a(n-9) - 2045028\,a(n-10) - 5829087\,a(n-11) + 9632462\,a(n-12) + 14361840\,a(n-13) - 27309714\,a(n-14) - 21401634\,a(n-15) + 48701640\,a(n-16) + 18436650\,a(n-17) - 56085498\,a(n-18) - 7403128\,a(n-19) + 42308694\,a(n-20) - 1027934\,a(n-21) - 20976128\,a(n-22) + 2642614\,a(n-23) + 6772235\,a(n-24) - 1305140\,a(n-25) - 1387863\,a(n-26) + 326941\,a(n-27) + 171884\,a(n-28) - 44771\,a(n-29) - 11807\,a(n-30) + 3188\,a(n-31) + 387\,a(n-32) - 103\,a(n-33) - 4\,a(n-34) + a(n-35)
\]
for every $n>35$.
\end{theorem}

\begin{proof}
The vector $w=q(M)\tau$ was formed by $35$ matrix--vector products in exact integer
arithmetic and $\iota^{\!\top}M^{j}w$ evaluated for $j=0,1,\dots$, again exactly; those
integers vanish from the index corresponding to $n=35+1$ onwards and the run continued
until the working vector was identically zero, which settles every larger $j$ at once.
\end{proof}

\section{Verification}

Four checks, each independent of the algebra above.

First, the model reproduces all $16$ terms the entry publishes, exactly, in integer
arithmetic.

Second, the count was recomputed for the smallest arrays straight from the definition: every
array of the stated width and a few rows was written out cell by cell over $\{0,1,2,3,4\}$, each of
its $2\times2$ subblocks was read round its perimeter, and the condition of Section 2 was
applied as stated. No window, no state and no recurrence is involved in that computation, and
the published terms came back.

Third, the conjectured recurrence was evaluated directly on the published terms, with no
matrices involved, and holds wherever the proved range and the data overlap.

Fourth, the annihilation test of Lemma~\ref{lem:crit} was carried out in exact integer
arithmetic on the whole vector rather than on a sampled prefix.

\begin{thebibliography}{9}
\bibitem{oeis} The OEIS Foundation, \emph{The On-Line Encyclopedia of Integer
Sequences}, \url{https://oeis.org}, 2026. Sequence A183716.
\bibitem{stanley} R.~P.~Stanley, \emph{Enumerative Combinatorics}, Volume 1, 2nd ed.,
Cambridge University Press, 2011. (Transfer-matrix method, Section 4.7.)
\bibitem{horn} R.~A.~Horn and C.~R.~Johnson, \emph{Matrix Analysis}, 2nd ed., Cambridge
University Press, 2013.
\end{thebibliography}

\end{document}
